\section{Cross-layer questions that organize the analysis}\label{sec:C01:07}
The first question is whether retained information remains usable. A memory system can archive every event and still fail to retrieve a rare instruction before an irreversible action. Retrieval can succeed while the policy fails to recognize implicit relevance. A summary can preserve a fact while dropping its temporal qualifier. These are different failures, requiring separate measurements of retention, access, interpretation, influence, and execution. Treating all of them as recall produces misleading explanations.

The second question is whether the system acts under the history it is trained to assume. Context compaction changes conditioning. Training reconstructed trajectories that expose information unavailable during rollout can create leakage or support mismatch. MemoryWalker explicitly studies this problem, while Context Language Models make context editing part of the policy interface. Their relationship is a cross-layer issue: an editing policy, a learning objective, and a serving cache must agree about which history produced each action~\cite{j2026memorywalker,shao2026context}.

The third question concerns improvement. A system that accumulates successful procedures may become better on recurring tasks while growing less robust to novel environments. An evolving harness can overfit a finite selection set. An evaluator can drift as the executor changes. Recursive improvement requires distinguishing the quality of the executor from the quality of the process that proposes, selects, and validates future changes. RRSI and Harness\textsuperscript{2} provide concrete but bounded mechanisms to analyze; neither licenses a general claim of indefinite acceleration~\cite{rrsi2026,corpus261014224}. The additional adaptation question is whether a revision improves its available control problem or merely fits the current policy; a frozen-policy ranking cannot settle that distinction. Section~\ref{sec:C01:interaction} specifies the corresponding comparison boundary.

The fourth question is operational: which guarantees survive interruption, concurrency, migration, and partial failure? A correct answer in a clean benchmark environment may rely on assumptions absent in deployment. Duplicate actions, inconsistent checkpoints, stale indexes, revoked permissions, and unobserved external effects connect systems engineering directly to model behavior. Reliable autonomy is an end-to-end property of this coupled process. The subsequent chapters develop the mathematical, computational, and empirical machinery needed to test it.
